\chapter{Drugs in Sport}
\index{Drugs in Sport}

\section*{Definition and Overview}
Drugs in sport refers to the use of substances or methods by athletes to enhance performance, alter body composition, or mask other drug use. Doping is prohibited by sporting organisations worldwide because it is unfair to clean athletes, dangerous to health, and contrary to the spirit of sport. The World Anti-Doping Agency (WADA) maintains the Prohibited List, which details the substances and methods banned in and out of competition. This chapter outlines the main categories of performance-enhancing drugs, the health risks they carry, and the rules that govern their use, including the therapeutic use exemption process for athletes who need banned medicines for legitimate medical reasons.

\section*{Epidemiology}
Doping is difficult to measure precisely because it is concealed. Population surveys suggest that a small proportion of elite athletes admit to doping, while recreational gym users use substances such as anabolic steroids and stimulants at higher rates than is often assumed. The prevalence of doping varies by sport, level of competition, and country. Anti-doping testing programmes now carry out hundreds of thousands of samples each year worldwide, and the proportion of adverse analytical findings is small but persistent, indicating that a minority of athletes continue to attempt to gain an unfair advantage.

\section*{Aetiology and Pathophysiology}
Doping arises from the powerful incentives of sporting success, including prize money, sponsorship, and national pride, combined with the belief that enhancement is necessary to compete. The substances used act on the body in different ways.
\begin{itemize}
    \item \textbf{Anabolic-androgenic steroids:} synthetic versions of testosterone that increase muscle mass and strength by stimulating protein synthesis; also used in stacks with other hormones
    \item \textbf{Stimulants:} drugs such as amphetamines and ephedrine that increase alertness, aggression, and endurance by activating the sympathetic nervous system
    \item \textbf{Erythropoietin (EPO) and blood doping:} increase the oxygen-carrying capacity of blood, improving endurance performance
    \item \textbf{Growth hormone and insulin-like growth factors:} promote muscle growth and fat loss
    \item \textbf{Beta-blockers:} paradoxically banned in some sports such as shooting and archery because they steady the hands by slowing the heart
    \item \textbf{Diuretics and masking agents:} used to flush other drugs from the body or rapidly reduce weight in weight-category sports
    \item \textbf{Gene doping and novel methods:} manipulating genetic material or cells to enhance performance, an emerging and prohibited frontier
\end{itemize}

\section*{Clinical Features}
The health effects of doping drugs are wide-ranging and often serious.
\begin{itemize}
    \item Acne, hair loss, and oily skin with anabolic steroid use
    \item Mood changes, aggression, irritability, and depression, sometimes called "roid rage"
    \item Testicular shrinkage, reduced sperm count, and infertility in men; menstrual irregularities and masculinising effects in women
    \item Increased risk of heart attack, stroke, and high blood pressure
    \item Liver damage, including tumours, with long-term steroid use
    \item Psychological dependence on performance-enhancing drugs and body image preoccupation
    \item Infections from unsafe injecting, including abscesses, hepatitis, and HIV
\end{itemize}

\section*{Differential Diagnosis}
Athletes may use substances for reasons other than performance enhancement.
\begin{itemize}
    \item Social or recreational drug use outside the sporting context
    \item Misuse of prescribed medicines for pain, anxiety, or sleep, without intention to enhance performance
    \item Legitimate medical treatment requiring a banned substance, covered by a therapeutic use exemption
    \item Use of supplements unknowingly contaminated with prohibited substances
    \item Body dysmorphic disorder and muscle dysmorphia driving appearance-focused steroid use
\end{itemize}

\section*{When to Seek Medical Care}
Athletes who are using or considering using performance-enhancing drugs should speak with a GP or sports medicine physician, who can discuss the health risks and provide non-doping alternatives to improve performance. Anyone experiencing concerning symptoms such as chest pain, severe mood disturbance, or depression related to substance use should seek medical care promptly.

\textbf{Contact your local emergency services for:}
\begin{itemize}
    \item Chest pain, palpitations, or collapse during or after training
    \item Severe agitation, paranoia, hallucinations, or aggressive behaviour
    \item Signs of a stroke, including sudden weakness, facial drooping, or difficulty speaking
    \item Thoughts of self-harm or suicide
\end{itemize}

\section*{Diagnosis and Investigations}
Detection of doping relies on the anti-doping testing system.
\begin{itemize}
    \item \textbf{Urine and blood testing:} samples collected in and out of competition are analysed in accredited laboratories for prohibited substances
    \item \textbf{Athlete biological passport:} long-term monitoring of selected biological markers detects the effects of doping rather than the substance itself
    \item \textbf{Therapeutic use exemptions:} athletes with a genuine medical need for a prohibited medicine can apply for approval before using it
    \item \textbf{Medical assessment:} athletes with health concerns should have blood pressure, liver function, hormone levels, and mental health reviewed
\end{itemize}

\section*{Management}
Management of doping in sport operates at the level of prevention, detection, and treatment.
\begin{itemize}
    \item \textbf{Education:} athletes, coaches, and support staff are educated about the Prohibited List, the risks of doping, and the rules
    \item \textbf{Testing programmes:} in-competition and out-of-competition testing, with sanctions for positive results
    \item \textbf{Healthy performance strategies:} evidence-based nutrition, strength and conditioning, recovery, and sports psychology offer safe performance gains
    \item \textbf{Support for athletes:} athletes who test positive are offered education, counselling, and, where appropriate, treatment for substance dependence
    \item \textbf{Monitoring:} doctors monitoring athletes who have used anabolic steroids for muscle wasting or low testosterone should follow established protocols
\end{itemize}

\section*{Complications}
\begin{itemize}
    \item Serious cardiovascular, liver, hormonal, and psychiatric complications from doping drugs
    \item Bans, fines, loss of medals, and damage to reputation and career
    \item Psychological dependence and body image disorders
    \item Blood-borne infections from injecting practices
    \item Unfair competition, which undermines the integrity of sport and the trust of fans
\end{itemize}

\section*{Prognosis and Outcomes}
The prognosis for doping-related health problems depends on the substance, duration of use, and how early problems are detected. Many steroid-related effects, such as reduced fertility and mood disturbance, improve when use stops, although some cardiovascular and liver damage may persist. Sport itself is cleaner and fairer when anti-doping programmes are strong, and the majority of athletes compete honestly. Athletes who are caught face serious consequences, but support is available to help them move forward.

\section*{Prevention and Self-Care}
\begin{itemize}
    \item Compete clean: success built on hard training and good preparation is sustainable and satisfying
    \item Know the rules: check the WADA Prohibited List before using any medicine or supplement
    \item Use the global drug reference database to check whether a medication is permitted
    \item Beware of supplements: many are unregulated and some are contaminated with prohibited substances
    \item Never share needles or equipment
    \item Seek advice from a sports medicine professional rather than relying on gym folklore
    \item Talk openly with young athletes about the risks and the value of fairness in sport
\end{itemize}

\section*{Sources and Further Reading}
The following reputable sources provide further information for healthcare students and patients seeking reliable, current advice about drugs in sport.
\begin{itemize}
    \item World Anti-Doping Agency. \textit{The Prohibited List and anti-doping rules.} https://www.wada-ama.org/
    \item Sport Integrity Australia. \textit{Drug testing and anti-doping information for athletes.} https://www.sportintegrity.gov.au/
    \item NHS. \textit{Steroids: side effects and health risks.} https://www.nhs.uk/conditions/anabolic-steroid-misuse/
    \item Mayo Clinic. \textit{Performance-enhancing drugs: know the risks.} https://www.mayoclinic.org/healthy-lifestyle/fitness/
    \item MedlinePlus. \textit{Performance-enhancing drugs.} https://medlineplus.gov/performanceenhancingdrugs.html
    \item National Institute on Drug Abuse. \textit{Anabolic steroids and athletic performance.} https://nida.nih.gov/research-topics/anabolic-steroids
\end{itemize}
